\documentclass[11pt,a4paper,oneside]{book}
\usepackage{mathscommun}
\setniveau{Seconde}
\renewcommand{\ReferenceProgramme}{Programme français de seconde GT -- rentrée 2026}
\renewcommand{\EditionDocument}{Édition du 18 octobre 2026}
\hypersetup{pdftitle={Mathématiques de seconde - Guide, programme et révisions}}
\begin{document}
\titrecollection{Guide, programme et révisions}{Organiser le travail et vérifier les acquis.}{Table de couverture du programme de seconde générale et technologique, preuves de référence, progression, rappels et soixante automatismes corrigés. Les approfondissements préparent progressivement un parcours vers les mathématiques universitaires.}
\frontmatter\tableofcontents\mainmatter
\chapter{Périmètre et sources}
\section{Le programme applicable}
Ce volume est aligné sur le programme de mathématiques de seconde générale et technologique publié au Bulletin officiel n\textsuperscript{o} 14 du 2 avril 2026, arrêté du 26 février 2026, NOR MENE2602914A. Il entre en application à la rentrée 2026--2027 et remplace l'annexe de 2019 pour cette classe.

Il comprend notamment les tableaux croisés, les fréquences et probabilités conditionnelles, les séries regroupées en classes et la fonction valeur absolue. Une table fondée uniquement sur le texte de 2019 manquerait donc des attendus de ce périmètre.
\begin{enumerate}
\item Ministère de l'Éducation nationale, \href{https://www.education.gouv.fr/bo/2026/Hebdo14/MENE2602914A}{arrêté, calendrier et annexe officielle du programme de seconde GT}. L'annexe compte treize pages.
\item Éduscol, \href{https://eduscol.education.gouv.fr/5817/programmes-et-ressources-en-mathematiques-voie-gt}{programmes et ressources de mathématiques pour la voie générale et technologique}.
\item STATEC, \href{https://statistiques.public.lu/dam-assets/catalogue-publications/en-chiffres/2025/population-chiffres-2025.pdf}{La démographie luxembourgeoise en chiffres, édition 2025}, tableau 2, p.~11 : effectifs au 1\textsuperscript{er} janvier 2025, utilisés au chapitre 20.
\end{enumerate}
Sources consultées le 18 octobre 2026. Les données cantonales sont un instantané daté ; le dossier de données indique la transcription et les contrôles. Les autres contextes numériques sont des exemples pédagogiques, fictifs lorsqu'aucune source n'est donnée.
\section{Le parcours vers le supérieur}
L'objectif final est d'acquérir des connaissances et une autonomie permettant d'étudier des mathématiques universitaires. Ce volume reste une étape de seconde. Il renforce l'exactitude du vocabulaire, la maîtrise du calcul, les preuves, les contre-exemples, la modélisation et l'interprétation des données.

La suite envisagée est la spécialité de première, puis la spécialité de terminale avec les compléments pertinents, notamment ceux de mathématiques expertes, avant des volumes universitaires d'analyse, d'algèbre linéaire et de probabilités. Les programmes applicables devront être vérifiés au moment de la rédaction. Ce parcours est une orientation pédagogique, pas une équivalence de diplôme.
\section{Les quatre documents}
\begin{center}\begin{tabularx}{\textwidth}{lX}\toprule
Document & Utilisation\\\midrule
Cours & 24 chapitres, définitions, méthodes, preuves et figures.\\
Exercices & 144 exercices, six par chapitre, et six synthèses.\\
Corrigés & Même numérotation, démarches et contrôles expliqués.\\
Guide & Couverture, fondements, progression, rappels et 60 automatismes.\\\bottomrule
\end{tabularx}\end{center}
Les chapitres 1 et 23 servent de références transversales ; ils ne doivent pas monopoliser plusieurs semaines avant les autres thèmes. Le chapitre 24 est entièrement facultatif. Des compléments ponctuels sont aussi signalés dans le cours.
\chapter{Table de couverture du programme}
Cette table croise les rubriques et capacités de l'annexe officielle. Les références « A » désignent les automatismes du présent guide. Les numéros de type 19.3 désignent les exercices, dont les corrigés portent les mêmes numéros. Une notion peut être travaillée dans plusieurs chapitres.
\section{Vocabulaire ensembliste et logique}
\small
\begin{longtable}{p{.48\textwidth}p{.14\textwidth}p{.27\textwidth}}
\toprule Notions et capacités & Cours & Entraînement\\\midrule
\endfirsthead\toprule Notions et capacités & Cours & Entraînement\\\midrule\endhead
Éléments, inclusion, ensemble vide ; union, intersection et complémentaire ; cardinal & 1 & 1.1, 1.3, 1.5\\
Couples et produit cartésien ; dénombrement & 1, 21 & 1.2 ; 21.2, 21.6\\
Propositions, variables, ET, OU, négation simple ; contre-exemples & 1, 12, 23 & 1.3--1.4 ; 12.4\\
Implication, réciproque, contraposée, équivalence ; quantification en français & 1, 2, 6 & 1.4, 1.6 ; 2.6\\
Preuves directes, par cas, par l'absurde & 1--3, 9 & 2.3, 2.6 ; 3.5 ; 9.4\\
\bottomrule\end{longtable}
\normalsize
\section{Nombres et calculs, algèbre}
\small
\begin{longtable}{p{.48\textwidth}p{.14\textwidth}p{.27\textwidth}}
\toprule Notions et capacités & Cours & Entraînement\\\midrule
\endfirsthead\toprule Notions et capacités & Cours & Entraînement\\\midrule\endhead
Entiers naturels et relatifs ; multiples, diviseurs et parité & 2 & 2.1--2.2, 2.4--2.6\\
Écriture fractionnaire irréductible ; utilisation des nombres premiers & 2 & 2.3 ; A1--A4\\
Ensembles de nombres, décimaux, rationnels, irrationnels & 3 & 3.1, 3.5\\
Droite réelle, intervalles, réunion et intersection ; distance et valeur absolue & 3, 14 & 3.2--3.3 ; 14.1\\
Encadrements décimaux, arrondis et chiffres significatifs & 3, 23 & 3.4, 3.6 ; 23.3\\
Puissances entières, écriture scientifique et racines carrées & 4 & 4.1--4.6\\
Expressions algébriques, identités, développement, factorisation ; formes adaptées & 5, 14, 16 & 5.1--5.2, 5.5--5.6\\
Quotients, valeurs interdites, opérations sur fractions rationnelles & 5--6 & 5.3 ; 6.3\\
Relations d'ordre : addition, multiplication de signe connu ; comparaison & 6 & 6.4--6.5 ; A2, A17\\
Équations affines, produits nuls, carrés, quotients égaux à une constante & 6, 14--15 & 6.1--6.3 ; 14.4 ; 15.1--15.2\\
Inéquations affines, produits et quotients ; tableaux de signes & 6, 13, 15 & 6.4--6.6 ; 15.1\\
Isoler une variable ; formules physiques et géométriques & 5 & 5.4 ; A18--A20\\
Algorithmique : test de divisibilité, dernier multiple, seuil de puissance, balayage d'une racine & 2--4, 23 & 2.2 ; 4.5 ; 23.2--23.3\\
\bottomrule\end{longtable}
\normalsize
\section{Géométrie}
\small
\begin{longtable}{p{.48\textwidth}p{.14\textwidth}p{.27\textwidth}}
\toprule Notions et capacités & Cours & Entraînement\\\midrule
\endfirsthead\toprule Notions et capacités & Cours & Entraînement\\\midrule\endhead
Repère orthonormé, coordonnées ; milieu et distance & 7 & 7.1--7.6\\
Vecteur, translation, égalité, vecteur nul ; norme & 8 & 8.1, 8.6\\
Addition, Chasles, multiplication par un réel ; combinaison linéaire & 8 & 8.2--8.5\\
Colinéarité, coordonnées proportionnelles et déterminant & 9 & 9.1, 9.4\\
Alignement et parallélisme par plusieurs méthodes & 8--11 & 9.2--9.3, 9.6 ; 11.1\\
Décomposition dans une base non colinéaire & 9 & 9.5\\
Droites : vecteur directeur, équations cartésienne et réduite, pente & 10, 13 & 10.1--10.3\\
Construire une droite par deux points, point et directeur, point et pente & 10 & 10.1--10.3 ; A33\\
Position relative, parallélisme, intersection ; systèmes linéaires & 10 & 10.4--10.5 ; 9.6\\
Calculer et démontrer dans des configurations ; choix d'une méthode & 7--11 & 7.3--7.6 ; synthèse 2\\
Algorithmique : alignement, équation d'une droite & 9--10, 23 & 23 ; module Python\\
\bottomrule\end{longtable}
\normalsize
\section{Fonctions}
\small
\begin{longtable}{p{.48\textwidth}p{.14\textwidth}p{.27\textwidth}}
\toprule Notions et capacités & Cours & Entraînement\\\midrule
\endfirsthead\toprule Notions et capacités & Cours & Entraînement\\\midrule\endhead
Variable, domaine, image, antécédent ; notation f(x) et u(n) & 12 & 12.1--12.6\\
Formule, tableau, courbe ; appartenance d'un point ; modèle discret ou continu & 12 & 12.3--12.6 ; synthèse 3\\
Fonctions affines, pente comme taux de variation et interpolation & 13 & 13.1--13.6\\
Valeur absolue et carré : définition, graphes, signes et ordre & 14 & 14.1--14.6\\
Inverse, racine carrée et cube : domaines, graphes et variations & 15 & 15.1--15.6\\
Résoudre f(x)=k, f(x)<k, f(x)=g(x), f(x)<g(x) algébriquement ou graphiquement & 6, 12--15 & 6.4--6.5 ; 12.1 ; 13.3 ; 15.2\\
Croissance, décroissance, monotonie ; tableaux de variations & 13--16 & 13.4 ; 16.1\\
Maximum et minimum sur un domaine, y compris leur existence ; optimisation & 14, 16 & 14.5 ; 16.2 ; synthèse 1\\
Algorithmique : balayage, dichotomie liée à un extremum, longueur polygonale & 16, 23 & 16.3--16.6 ; 23.6\\
\bottomrule\end{longtable}
\normalsize
\section{Statistiques et probabilités}
\small
\begin{longtable}{p{.48\textwidth}p{.14\textwidth}p{.27\textwidth}}
\toprule Notions et capacités & Cours & Entraînement\\\midrule
\endfirsthead\toprule Notions et capacités & Cours & Entraînement\\\midrule\endhead
Proportions et proportion de proportion ; changement de population de référence & 17, 20 & 17.1--17.2 ; 20.2\\
Variation absolue, relative, coefficient multiplicateur ; points de pourcentage & 17 & 17.3--17.5\\
Évolutions successives et réciproques ; calculer un taux ou une valeur & 17 & 17.3--17.4 ; A37--A40\\
Moyenne pondérée et linéarité ; médiane ; ajout ou retrait d'une valeur & 18 & 18.1--18.4\\
Quartiles, étendue, écart interquartile ; diagramme en boîte & 18 & 18.1, 18.6 ; synthèse 4\\
Variance descriptive et écart-type ; effets d'une transformation affine & 19 & 19.1--19.2\\
Classes de même amplitude : histogramme et polygone des fréquences cumulées & 19 & 19.3\\
Moyenne par moyennes de classes ; estimation par centres et hypothèse uniforme & 19 & 19.3--19.4\\
Classe médiane et interpolation ; comparer deux distributions & 19 & 19.3, 19.5 ; synthèse 4\\
Variables qualitatives nominales ou ordinales ; tableaux croisés & 20 & 20.1, 20.3\\
Effectifs, marges, fréquences conjointes, marginales et conditionnelles & 20 & 20.1--20.2, 20.4--20.6\\
Données réelles et filtres ET, OU, NON ; tableau à partir de listes ou fichier & 20, 23 & 20.4--20.5 ; synthèse 6 ; CSV et Python\\
Univers fini, loi de probabilité, équiprobabilité, dénombrement & 21 & 21.1--21.2, 21.4, 21.6\\
Événements, contraire, union et intersection ; additivité & 1, 21 & 21.3 ; A49--A52\\
Répétitions indépendantes, fluctuation ; fréquences et probabilités & 21, 23 & 21.5 ; 23.5\\
Probabilité conditionnelle ; tableau et arbre pondéré ; inversion du conditionnement & 22 & 22.1--22.3, 22.6\\
Chemins disjoints et règle du produit ; calculs contextualisés & 22 & 22.2 ; synthèse 5\\
Tests : faux positifs, faux négatifs, sensibilité et spécificité & 22 & 22.3--22.4\\
\bottomrule\end{longtable}
\normalsize
\section{Algorithmique et automatismes transversaux}
\small
\begin{longtable}{p{.48\textwidth}p{.14\textwidth}p{.27\textwidth}}
\toprule Notions et capacités & Cours & Entraînement\\\midrule
\endfirsthead\toprule Notions et capacités & Cours & Entraînement\\\midrule\endhead
Types entier, flottant, booléen, chaîne ; affectation et séquence & 23 & 23.1 ; A53--A54\\
Conditions, boucles for et while ; lecture, modification, écriture & 2, 4, 23 & 23.1--23.3 ; A55\\
Fonctions à un ou plusieurs arguments, appels ; nombres aléatoires & 23 & 23.2, 23.5\\
Lire une fonction de moyenne et d'écart-type ; listes non exigées & 23 & 23.4\\
Calcul mental, fractions, puissances, écritures, plausibilité et ordre de grandeur & 2--6 & A1--A20\\
Conversions longueur, aire, volume, contenance, durée, vitesse et masse & Guide, rappels & A21--A27\\
Périmètres, aires, volumes usuels ; Pythagore, Thalès, trigonométrie & 7, 11 ; guide & A28--A36 ; cours de 3e, ch. 12--17\\
Lire et exploiter les graphiques usuels, dont barres, secteurs, courbes et nuages & 12--13, 18--20 & 18.5 ; A41--A48, A56--A60\\
\bottomrule\end{longtable}
\normalsize

\section{Les six compétences}
Chercher : essayer plusieurs démarches et résoudre les synthèses.
Modéliser : choisir un domaine, un tarif, une distribution ou une loi de probabilité.
Représenter : passer des mots aux tableaux, figures, courbes et arbres.
Raisonner : prouver, réfuter et contrôler une réciproque.
Calculer : utiliser le mental, l'exact, l'approché et les outils.
Communiquer : rédiger des hypothèses, des étapes et une conclusion interprétée.
\chapter{Preuves, dépendances et approfondissements}
\section{Les démonstrations exemplaires du programme}
Le programme identifie les démonstrations suivantes ; elles figurent toutes dans le cours.
\begin{center}\begin{tabularx}{\textwidth}{Xl}\toprule
Résultat & Chapitre\\\midrule
Somme de deux multiples d'un entier & 2\\
Carré d'un entier impair & 2\\
Non-décimalité de $1/3$ & 3\\
Irrationalité de $\sqrt2$ & 3\\
Racine d'un produit de réels positifs ou nuls & 4\\
Colinéarité : proportionnalité des coordonnées et déterminant nul & 9\\
Équation cartésienne d'une droite & 10\\
Variation d'une fonction affine suivant le signe de sa pente & 13\\
Comparaison de x et $x^2$ pour $x\ge0$ & 14\\
Variations de la fonction carré & 14\\
Variations de la fonction inverse & 15\\\bottomrule
\end{tabularx}\end{center}
D'autres preuves sont développées : opérations sur fractions, identités, règles d'ordre, produit nul, distances et milieux, propriétés vectorielles, existence et unicité d'une décomposition dans une base, systèmes, variations de racine et cube, optimisation, coefficients d'évolution, linéarité de la moyenne, variance, règles de probabilité et correction d'algorithmes.
\section{Un socle admis explicite}
\begin{enumerate}
\item \textbf{Réels :} construction de $\R$, propriétés fondamentales de corps ordonné et propriété archimédienne. L'existence des racines carrées positives et des racines cubiques est admise ; leur unicité est justifiée par la monotonie.
\item \textbf{Arithmétique :} unicité de la décomposition en facteurs premiers. L'existence et plusieurs conséquences sont expliquées ; le lemme d'Euclide n'est pas reconstruit.
\item \textbf{Géométrie :} cadre euclidien, existence des parallèles et projections, translations et coordonnées. Pythagore et Thalès, déjà démontrés en troisième, sont mobilisés sans raisonnement circulaire.
\item \textbf{Mesures :} additivité et propriétés usuelles des aires et volumes ; passages à la limite nécessaires aux figures courbes. Le volume de troisième explique les formules de solides sous ces principes.
\item \textbf{Analyse :} continuité et valeurs intermédiaires des fonctions usuelles lorsqu'elles servent à une existence ; convergence des longueurs polygonales pour une courbe régulière, sans preuve générale à ce niveau.
\item \textbf{Probabilités :} choix du modèle et de ses probabilités élémentaires. Une simulation ne démontre pas la validité du modèle. Le rapprochement des fréquences est expliqué dans le tronc du cours et une preuve quantitative est proposée au chapitre 24 pour des épreuves de Bernoulli indépendantes.
\item \textbf{Récurrence :} principe sur les entiers admis dans le complément, puis appliqué avec initialisation et hérédité.
\end{enumerate}
Une définition n'a pas à être démontrée. Une convention, comme celle des quartiles ou des arrondis, doit être annoncée. Une formule d'estimation sous hypothèse uniforme ne devient pas une valeur exacte de la série brute.
\section{Les compléments facultatifs}
Sont signalés comme approfondissements : périodicité décimale des rationnels ; algorithme d'Euclide ; identités à trois termes et cube d'une somme ; inégalité arithmético-géométrique ; barycentres et concours de droites d'un triangle ; aire par sinus ; lieu parabolique et demi-plans ; symétrie des graphes carré/racine ; décomposition de variance ; indépendance formalisée, formule générale des probabilités totales et formule de Bayes ; quantificateurs symboliques, récurrence et loi faible des grands nombres.

La formule générale des probabilités totales relativement à une partition n'est pas exigée en seconde. Le calcul concret sur les chemins disjoints d'un arbre est, lui, travaillé dans le programme. De même, les symboles $\forall$ et $\exists$ ne font pas partie des attendus, contrairement aux phrases « pour tout » et « il existe ».
\section{Lire les dépendances d'une preuve}
La formule de distance s'appuie sur Pythagore ; elle ne sert pas ensuite à prétendre prouver Pythagore sans autre fondement. Les variations du carré sont établies par la différence de carrés ; celles de la racine en découlent. La règle du produit dans un arbre vient de la définition d'une probabilité conditionnelle. Les encadrements numériques sont justifiés par la monotonie et un invariant de boucle.
\chapter{Progression et méthode}
\section{Un parcours en spirale}
Les périodes ci-dessous sont indicatives. Les reprises sont espacées ; la logique et la programmation s'intègrent aux notions rencontrées.
\begin{center}\small\begin{tabularx}{\textwidth}{lXX}\toprule
Période & Notions principales & Reprises et outils\\\midrule
1 & Nombres, intervalles, vecteurs, proportions (2--4, 7--8, 17) & Vocabulaire logique ; tests de divisibilité ; automatismes\\
2 & Algèbre, équations, fonctions, droites (5--6, 9--10, 12--13) & Contre-exemples ; calcul exact ; tableaux et tracés\\
3 & Fonctions de référence, variations et configurations (11, 14--16) & Preuves ; dichotomie ; optimisation\\
4 & Statistiques, classes et tableaux croisés (18--20) & Proportions ; lecture d'un CSV ; calcul de dispersion\\
5 & Probabilités et conditionnement (21--22) & Simulation ; données et arbres ; synthèses\\
Toute l'année & Références 1 et 23 ; reprises des acquis & Rédaction orale et écrite, erreurs fréquentes\\
Après maîtrise & Compléments choisis et chapitre 24 & Autonomie et transition vers la première\\\bottomrule
\end{tabularx}\end{center}
Ce découpage n'impose pas de terminer un domaine avant d'en commencer un autre. Par exemple, un problème de droites réactive simultanément équations, vecteurs et fonctions affines.
\section{Travailler une preuve}
Énoncer d'abord le résultat sans regarder le texte. Entourer les hypothèses. Pour chaque étape, nommer la définition ou la propriété utilisée. Refaire ensuite la preuve sur une feuille blanche, puis chercher ce qui échoue si une hypothèse disparaît.

Un exercice réussi par imitation immédiate doit être repris à distance. Pour valider un chapitre, pouvoir donner une définition précise, résoudre un exercice nouveau, justifier une méthode et corriger une erreur typique.
\section{Utiliser les outils}
Le calcul formel propose et vérifie des transformations ; développer les deux membres reste une preuve accessible d'identité. Le grapheur aide à conjecturer un signe ou un extremum ; le tableau de signes et l'algèbre les justifient. La géométrie dynamique aide à explorer un lieu ou un concours ; un raisonnement général conclut.

Dans un tableur, mettre les effectifs en colonne B et les valeurs en A ; une colonne C contenant le produit A fois B reconstitue les sommes pondérées. Un cumul se construit en ajoutant chaque effectif au cumul précédent. Utiliser des références absolues pour un total fixe lorsque l'on recopie une formule. Toujours contrôler la somme des fréquences et les unités.

\chapter{Réactiver les acquis du collège}
Ces rappels sont à utiliser au moment d'un besoin, pas comme un bloc de révision préalable. Les preuves complètes déjà accessibles en troisième sont indiquées ; les formules ne sont donc pas isolées de leur justification.
\section{Unités et grandeurs}
Si une longueur est multipliée par k, l'aire d'une figure semblable est multipliée par $k^2$ et le volume d'un solide semblable par $k^3$. Pour un rectangle, cela vient de $(ka)(kb)=k^2ab$ ; pour un pavé, de $(ka)(kb)(kc)=k^3abc$. Les découpages et la mesure étendent ces règles aux figures usuelles.

Ainsi $1\unite m=100\unite{cm}$, $1\unite{m^2}=10\,000\unite{cm^2}$ et $1\unite{m^3}=1\,000\,000\unite{cm^3}$. Par définition, $1\unite L=1\unite{dm^3}$, donc $1\unite{mL}=1\unite{cm^3}$. Une heure compte $60$ minutes et $3600$ secondes ; une vitesse se convertit en convertissant son numérateur et son dénominateur : $1\unite{km/h}=1000/3600\unite{m/s}$. Pour les masses, $1\unite{kg}=1000\unite g$.

Une aire ne s'exprime pas en mètres et un volume ne s'exprime pas en mètres carrés. Une durée de $1\unite h\,30\unite{min}$ vaut $1{,}5\unite h$, pas $1{,}30\unite h$. Ces contrôles doivent accompagner les résultats numériques.
\section{Périmètres, aires et volumes}
\begin{center}\small\begin{tabularx}{\textwidth}{lXX}\toprule
Objet & Formule & Justification et renvoi\\\midrule
Polygone & Périmètre : somme des côtés & Définition d'une longueur de contour polygonal\\
Cercle & $P=2\pi r$ & Définition de $\pi$ et proportionnalité par similitude ; 3e, ch. 12\\
Rectangle & $A=ab$ & Pavage par unités d'aire ; extension aux réels\\
Triangle & $A=bh/2$ & Assemblage de deux triangles en parallélogramme ; 3e, ch. 12\\
Disque & $A=\pi r^2$ & Découpage en secteurs et passage à la limite ; 3e, ch. 12\\
Pavé & $V=abc$ & Empilement d'unités de volume ; 3e, ch. 17\\
Prisme, cylindre & $V=Bh$ & Sections de même aire B et hauteur h ; 3e, ch. 17\\
Pyramide, cône & $V=Bh/3$ & Découpage et principe de Cavalieri admis ; 3e, ch. 17\\
Boule & $V=4\pi r^3/3$ & Comparaison de sections avec cylindre et cônes ; 3e, ch. 17\\\bottomrule
\end{tabularx}\end{center}
La hauteur est perpendiculaire à la base. Les figures courbes demandent des principes de mesure et de passage à la limite qui ne sont pas reconstruits ici.
\section{Trois outils géométriques}
Dans un triangle ABC rectangle en A, $BC^2=AB^2+AC^2$ ; sa réciproque permet de prouver qu'un triangle est rectangle lorsque le carré du plus grand côté égale la somme des deux autres. La démonstration par découpage est donnée en troisième, chapitre 13.

Si A,M,B sont alignés, A,N,C sont alignés, les points sont dans la configuration de Thalès et $(MN)$ est parallèle à $(BC)$, les côtés correspondants sont proportionnels : $AM/AB=AN/AC=MN/BC$. Les dénominateurs sont non nuls et la configuration (emboîtée ou papillon) doit être identifiée. Pour la réciproque, vérifier aussi l'ordre des points sur les droites. Preuve par aires et réciproque : troisième, chapitre 14.

Pour un angle aigu $\theta$ d'un triangle rectangle,
$\cos\theta=\text{adjacent}/\text{hypoténuse}$,
$\sin\theta=\text{opposé}/\text{hypoténuse}$ et
$\tan\theta=\text{opposé}/\text{adjacent}$.
La similitude justifie que ces rapports dépendent uniquement de l'angle (troisième, chapitre 15). Diviser les deux premiers donne $\tan\theta=\sin\theta/\cos\theta$ ; diviser Pythagore par le carré de l'hypoténuse donne $\cos^2\theta+\sin^2\theta=1$. Régler la calculatrice en degrés pour les angles de ces exercices.
\chapter{Soixante automatismes corrigés}
Répondre sans calculatrice lorsque c'est possible. Mélanger plusieurs thèmes à chaque séance ; ces questions ne remplacent pas les exercices de recherche.
\section{Questions}
\begin{enumerate}
\item Calculer $\frac23+\frac16$.
\item Comparer $-\frac34$ et $-\frac23$.
\item Calculer $\frac35\div\frac9{10}$.
\item Rendre $84/126$ irréductible.
\item Calculer $2^{-3}\times2^5$.
\item Écrire $0{,}00072$ en notation scientifique.
\item Écrire $3/8$ en décimal et en pourcentage.
\item Donner un ordre de grandeur de $198\times51$.
\item Calculer $\sqrt{(-9)^2}$ et simplifier $\sqrt{75}$.
\item Donner les solutions de $|x-2|\le3$.
\item Développer $(x-4)^2$.
\item Factoriser $9x^2-25$.
\item Résoudre $5x-2=3x+8$.
\item Résoudre $x^2=7$.
\item Résoudre $6/x=3$ dans son domaine.
\item Résoudre $-2x+3>7$.
\item Pour $a,b>0$, que signifie $a/b<1$ ?
\item Dans $U=RI$, isoler R puis calculer pour $U=12$, $I=0{,}5$.
\item Un disque a une aire $49\pi\unite{cm^2}$. Donner son rayon.
\item Dans $2x+3y=12$, isoler y.
\item Convertir $2{,}7\unite{km}$ en mètres.
\item Convertir $3{,}2\unite{m^2}$ en $\unite{cm^2}$.
\item Convertir $0{,}004\unite{m^3}$ en litres.
\item Convertir $750\unite{mL}$ en $\unite{cm^3}$.
\item Écrire $1\unite h\,24\unite{min}$ en heures décimales et en secondes.
\item Convertir $72\unite{km/h}$ en $\unite{m/s}$.
\item Convertir $2{,}5\unite{kg}$ en grammes.
\item Donner le périmètre et l'aire d'un rectangle de côtés $3$ et $7$ cm.
\item Donner la circonférence et l'aire d'un disque de rayon $3$ cm.
\item Donner l'aire d'un triangle de base $8$ cm et de hauteur associée $5$ cm.
\item Donner le volume d'un pavé de dimensions $2,3,4$ cm et d'un prisme de base $12\unite{cm^2}$, hauteur $5$ cm.
\item Donner les volumes d'un cylindre de rayon $2$, hauteur $3$ ; d'un cône de mêmes dimensions ; d'une pyramide de base $9$, hauteur $4$ ; d'une boule de rayon $3$ (unités cohérentes).
\item Une droite passe par $(0;2)$ et a une pente $-3$. Donner son équation.
\item Un triangle rectangle a des côtés de l'angle droit $5$ et $12$. Donner l'hypoténuse.
\item Deux triangles en configuration de Thalès ont un rapport de réduction $2/3$. À un côté de $15$ du grand correspond quelle longueur du petit ?
\item Dans un triangle rectangle, un angle aigu vaut $30\degre$ et l'hypoténuse $10$. Donner le côté opposé et le côté adjacent.
\item Calculer $15\,\%$ de $240$.
\item Une baisse de $20\,\%$ suivie d'une hausse de $20\,\%$ conserve-t-elle la valeur initiale ?
\item Quelle hausse compense une baisse de $25\,\%$ ?
\item Un taux passe de $10\,\%$ à $12\,\%$. Donner l'écart en points et la variation relative.
\item Pour $f(x)=2x-3$, calculer $f(-1)$ et l'antécédent de $7$.
\item Le point $(2;5)$ appartient-il à la courbe de $g(x)=x^2+1$ ?
\item Donner le minimum de $(x-3)^2+2$ sur $\R$.
\item Comparer $1/3$ et $1/5$ à l'aide d'une variation.
\item Donner moyenne et médiane de $1,2,6,7$.
\item Donner $Q_1,Q_3$ et l'étendue de $1,2,3,4,5,6,7,8$.
\item Une série a une moyenne $8$ et un écart-type $2$. On ajoute $5$ à chaque valeur. Quels sont les nouveaux indicateurs ?
\item Une catégorie compte $30$ individus sur $120$. Quel angle représente-t-elle dans un disque puis dans un demi-disque ?
\item Si $P(A)=0{,}35$, calculer $P(\overline A)$.
\item Sur un dé équilibré, calculer la probabilité d'un résultat strictement supérieur à $4$.
\item Si $P(A)=0{,}4$, $P(B)=0{,}3$ et $P(A\cap B)=0{,}1$, calculer $P(A\cup B)$.
\item Si $P(A)=0{,}5$ et $P_A(B)=0{,}2$, calculer $P(A\cap B)$.
\item Après \texttt{x = 3} puis \texttt{x = x + 2}, que vaut x ?
\item Donner les types Python de \texttt{4}, \texttt{4.0}, \texttt{True} et \texttt{"4"}.
\item Quelles valeurs parcourt \texttt{range(2, 6)} ?
\item Une courbe passe par $(3;8)$. Traduire cette information en notation fonctionnelle.
\item Dans un diagramme en barres, deux catégories ont des effectifs $12$ et $18$. Quel rapport de hauteurs faut-il si l'origine est zéro et l'échelle commune ?
\item Un nuage de points montre une association entre deux mesures. Cela prouve-t-il un lien causal ?
\item Un cumul de fréquences vaut $0{,}2$ à $10$ et $0{,}6$ à $20$. Estimer où il atteint $0{,}5$ si le segment est linéaire.
\item Une mesure de masse de $70\unite{kg}$ est saisie comme $70\,000\unite{kg}$. Quel contrôle aurait pu détecter l'erreur ?
\end{enumerate}
\section{Réponses et contrôles}
\begin{longtable}{r p{.84\textwidth}}
\toprule N\textsuperscript{o} & Réponse\\\midrule
\endfirsthead\toprule N\textsuperscript{o} & Réponse\\\midrule\endhead
1 & $\frac56$ en prenant le dénominateur $6$.\\
2 & $-\frac34<-\frac23$ car $-\frac9{12}<-\frac8{12}$.\\
3 & $\frac35\times\frac{10}{9}=\frac23$.\\
4 & $2/3$, en divisant par $42$.\\
5 & $2^2=4$.\\
6 & $7{,}2\times10^{-4}$.\\
7 & $0{,}375=37{,}5\,\%$.\\
8 & $200\times50=10\,000$ ; le produit exact est $10\,098$.\\
9 & $9$ et $5\sqrt3$.\\
10 & $[-1;5]$.\\
11 & $x^2-8x+16$.\\
12 & $(3x-5)(3x+5)$.\\
13 & $x=5$.\\
14 & $x=\pm\sqrt7$.\\
15 & $x\ne0$ puis $x=2$.\\
16 & $x<-2$. Le signe s'inverse à la division par $-2$.\\
17 & $a<b$, car on multiplie par $b>0$.\\
18 & $R=U/I$ si $I\ne0$ ; ici $R=24$ dans les unités cohérentes de la formule.\\
19 & $r=\sqrt{S/\pi}=7\unite{cm}$.\\
20 & $y=4-\frac23x$.\\
21 & $2700\unite m$.\\
22 & $32\,000\unite{cm^2}$ : le facteur $100$ est élevé au carré.\\
23 & $4\unite L$, car $1\unite{m^3}=1000\unite L$.\\
24 & $750\unite{cm^3}$, puisque $1\unite{mL}=1\unite{cm^3}$.\\
25 & $1{,}4\unite h$ et $5040\unite s$.\\
26 & $72\times1000/3600=20\unite{m/s}$.\\
27 & $2500\unite g$.\\
28 & $20\unite{cm}$ et $21\unite{cm^2}$.\\
29 & $6\pi\unite{cm}$ et $9\pi\unite{cm^2}$.\\
30 & $8\times5/2=20\unite{cm^2}$.\\
31 & $24\unite{cm^3}$ et $60\unite{cm^3}$.\\
32 & $12\pi$, $4\pi$, $12$ et $36\pi$ unités de volume.\\
33 & $y=-3x+2$.\\
34 & $\sqrt{25+144}=13$.\\
35 & $15\times2/3=10$, après vérification du parallélisme et des correspondances.\\
36 & $10\sin30\degre=5$ et $10\cos30\degre=5\sqrt3$.\\
37 & $36$.\\
38 & Non : $0{,}8\times1{,}2=0{,}96$, soit une baisse globale de $4\,\%$.\\
39 & $1/0{,}75-1=1/3$, soit environ $33{,}33\,\%$.\\
40 & $2$ points et $20\,\%$ d'augmentation relative.\\
41 & $f(-1)=-5$ ; $2x-3=7$ donne $x=5$.\\
42 & Oui : $g(2)=5$.\\
43 & $2$, atteint en $x=3$.\\
44 & $1/3>1/5$ car l'inverse est décroissant sur $]0;+\infty[$.\\
45 & Moyenne $4$ ; médiane $(2+6)/2=4$.\\
46 & $Q_1=2$, $Q_3=6$, étendue $7$ (convention des rangs du cours).\\
47 & Moyenne $13$, écart-type $2$.\\
48 & $90\degre$ puis $45\degre$.\\
49 & $0{,}65$.\\
50 & $2/6=1/3$.\\
51 & $0{,}4+0{,}3-0{,}1=0{,}6$.\\
52 & $0{,}5\times0{,}2=0{,}1$.\\
53 & $5$ : l'affectation remplace l'ancienne valeur.\\
54 & \texttt{int}, \texttt{float}, \texttt{bool}, \texttt{str}.\\
55 & $2,3,4,5$ ; la borne finale est exclue.\\
56 & $f(3)=8$ : $8$ est l'image de $3$, et $3$ est un antécédent de $8$.\\
57 & $12/18=2/3$. Une origine tronquée peut déformer la comparaison visuelle.\\
58 & Non. Une variable tierce ou le protocole de collecte peut expliquer l'association.\\
59 & $10+10(0{,}5-0{,}2)/(0{,}6-0{,}2)=17{,}5$. L'hypothèse uniforme est nécessaire à cette interpolation.\\
60 & Comparer l'ordre de grandeur au contexte et vérifier l'unité : $70\,000$ grammes correspondent à $70$ kilogrammes.\\
\bottomrule\end{longtable}
\chapter{Sources, fichiers et vérification}
Le style unique \texttt{mathscommun.sty}, à la racine, est partagé avec la troisième. Les quatre maîtres de \texttt{seconde/} assemblent des fichiers par chapitre. Les figures utilisent TikZ/PGFPlots et restent vectorielles.

Le dossier \texttt{python/} contient des fonctions exécutables, des contrôles d'entrée et des vérifications portant sur des résultats mathématiques : encadrements, géométrie, seuils exacts, statistiques et marges du fichier cantonal. Les listes et la lecture de CSV sont des supports pratiques ; leur maîtrise générale n'est pas exigée par le programme de seconde.

Depuis la racine du dépôt :
\begin{lstlisting}
python3 compiler.py
python3 compiler.py --niveau seconde
python3 compiler.py --niveau seconde --chapitre 9
python3 seconde/python/verifier.py
\end{lstlisting}
Installer Python 3 et pdfLaTeX avec les extensions du style commun. Le script de compilation effectue trois passages et contrôle les erreurs, les débordements et les références. Les fichiers auxiliaires vont dans \texttt{build/} ; les PDF publiés dans \texttt{pdf/seconde/}. Un chapitre seul reste dans \texttt{build/}.

Toute évolution passe par une pull request : sources, PDF, calculs et table de couverture doivent rester cohérents. Une compilation réussie ne garantit pas à elle seule la validité mathématique ni la qualité visuelle ; ces deux vérifications restent nécessaires.
\end{document}
